\documentclass[11pt]{article}

\usepackage{amsmath,amssymb}
\usepackage{xurl}
\usepackage[hidelinks]{hyperref}
\setlength{\emergencystretch}{2em}

% Shared commands for notes in docs/paper-gaps/.
%
% Two halves.  The link commands below are project identity: OWNER/REPO is the
% placeholder that scripts/bootstrap_project.py replaces with this project's
% GitHub slug and Pages site.  Everything after them is NOTATION, and notation
% belongs to the source paper: the definitions here are an example set, and a
% project replaces them with the macros of its own paper's style file.  The one
% rule that never changes: a command defined here must mean exactly what the
% source paper's macro means, or a note silently says something else.

\newcommand{\Lean}{\textsc{Lean}}
\newcommand{\leanid}[1]{\textsf{#1}}
\newcommand{\ghissue}[1]{\href{https://github.com/OWNER/REPO/issues/#1}{GitHub issue \##1}}
\newcommand{\ghpr}[1]{\href{https://github.com/OWNER/REPO/pull/#1}{PR \##1}}
% Local issue tracker (issues/NNNN-*.md in this repository); no remote URL.
\newcommand{\localissue}[1]{issue \##1 (\path{issues/})}
\newcommand{\doclink}[1]{\href{https://OWNER.github.io/REPO/docs/#1.html}{\path{#1}}}

% --- Notation: replace with the source paper's own macros --------------------
% \F, \N, \C, \R, \tr, \ot, \eps, \E, \bx, \bu, \bv, \by

\newcommand{\F}{\mathbb{F}}
\newcommand{\N}{\mathbb{N}}
\newcommand{\C}{\mathbb{C}}
\newcommand{\R}{\mathbb{R}}
\newcommand{\tr}{\mathrm{tr}}
\newcommand{\eps}{\epsilon}
\newcommand{\ot}{\otimes}
\newcommand{\E}{\mathop{\bf E\/}}

% bold random variables
\newcommand{\bu}{\boldsymbol{u}}
\newcommand{\bv}{\boldsymbol{v}}
\newcommand{\bx}{{\boldsymbol{x}}}
\newcommand{\by}{\boldsymbol{y}}

% T1 so literal < and > in \texttt placeholders typeset as themselves
\usepackage[T1]{fontenc}

% bra-ket
\usepackage{braket}

% --- Example structured spaces ---
\newcommand{\polyfunc}[3]{\mathcal{P}(#1, #2, #3)}
\newcommand{\polymeas}[3]{\mathrm{PolyMeas}(#1, #2, #3)}
\newcommand{\polysub}[3]{\mathrm{PolySub}(#1, #2, #3)}

% Approximation relations (paper uses \approx_\eps and \simeq_\zeta)
\newcommand{\approxeps}{\approx_{\varepsilon}}
\newcommand{\simgeq}{\simeq_{\zeta}}


% Local notation matching the source paper (its style file is not shared here).
\newcommand{\val}{\operatorname{val}}
\newcommand{\game}{\mathfrak{G}}
\newcommand{\strategy}{\mathcal{S}}
\renewcommand{\ghissue}[1]{%
  \href{https://github.com/Dengnifer/MIPStarRE-A/issues/#1}{GitHub issue \##1}}
\renewcommand{\ghpr}[1]{%
  \href{https://github.com/Dengnifer/MIPStarRE-A/pull/#1}{PR \##1}}

\title{Attainment in the Symmetrization Lemma}
\author{}
\date{2026-09-12, realignment section added 2026-09-18}

\begin{document}
\maketitle

\section{At a glance}

\begin{itemize}
  \item \textbf{Difficulty: easy.}  The issue is a quantifier exchange between
  a supremum and an attained maximum in a preliminary lemma on nonlocal games.
  No analytic estimate is missing; the substance is deciding which statement
  the symmetrization proof supports and which statement the later arguments
  consume.
  \item \textbf{Estimated weight:} about \textbf{1--2 lemma interfaces} and on
  the order of \textbf{100--300 lines} of formal proof for the value-preserving
  symmetrization construction.  The attainment step itself should not be
  formalized as printed (see the Conclusion).
  \item \textbf{Mathlib/project split:} approximately \textbf{20\% Mathlib}
  (finite-dimensional inner product spaces, tensor products, block
  decompositions of $\C^2 \ot \C^d$) and \textbf{80\% project-local} (games,
  strategies, values, and the Naimark input).
  \item \textbf{Key mathematical inputs:} the completed projective form of
  Naimark dilation, \cite[Theorem~5.1]{Ji2020LowIndividualDegree}, and the
  definitions of symmetric games and tensor product values in the source
  \cite[Section~5]{Ji2020MIPStar}.  The note is audited against the
  proof-gap protocol in \ghissue{173}.
  \item \textbf{Formal domain obstruction:} the current Lean game structure
  also permits an empty answer alphabet.  Its strategy type is then empty,
  while the real conditional supremum defining its value is zero.  This gives
  a kernel-checked counterexample to the exact Lean statement at $\eps=1$;
  see the domain calculation below and \ghissue{524}.  It does not settle
  attainment for games with nonempty answer alphabets.
\end{itemize}

\section{Key theorem forms}

Three statements must be kept apart.  Throughout, $\game$ denotes a symmetric
two-player one-round game and $\eps \ge 0$ a real parameter; the notation is
introduced in the next section.

\begin{enumerate}
  \item \textbf{Source form (attainment form).}  If
  $\val^*(\game) = 1-\eps$, then there exists a symmetric and projective
  strategy $\strategy$ with $\val^*(\game, \strategy) \ge 1-\eps$.
  \item \textbf{Approximate form (what the source proof establishes).}  If
  $\val^*(\game) = 1-\eps$, then for every $\eps' > \eps$ there exists a
  symmetric and projective strategy $\strategy$ with
  $\val^*(\game, \strategy) \ge 1-\eps'$.
  \item \textbf{Given-strategy form (value-preserving symmetrization).}  If
  $\strategy'$ is any strategy for $\game$ with
  $\val^*(\game, \strategy') \ge 1-\eps$, then there exists a symmetric and
  projective strategy $\strategy$ with
  $\val^*(\game, \strategy) = \val^*(\game, \strategy') \ge 1-\eps$.
\end{enumerate}

The source proof consists exactly of a proof of the third form, applied to a
strategy supplied by the definition of the supremum; this yields the second
form.  The first form asserts strictly more than the second, and the passage
from the second to the first is the gap discussed in this note.  Every use of
the lemma, in the source and in this project, consumes only the second or the
third form.

\section{The source assertion}

A two-player one-round game
$\game = (\mathcal X, \mathcal A, \mu, D)$ in compact symmetric notation
consists of a finite question alphabet $\mathcal X$, a finite answer alphabet
$\mathcal A$, a symmetric question distribution $\mu$ on
$\mathcal X \times \mathcal X$, and a decision predicate $D$ treating the two
players symmetrically.  A tensor product strategy
$\strategy = (\ket\psi, A, B)$ consists of a unit vector $\ket\psi$ in a
tensor product $\mathcal H_A \ot \mathcal H_B$ of \emph{finite-dimensional}
Hilbert spaces together with POVMs $\{A^x_a\}$ and $\{B^x_a\}$ for the two
players; its value $\val^*(\game, \strategy)$ is the acceptance probability,
and the value of the game is the supremum
\begin{equation*}
  \val^*(\game) = \sup_{\strategy} \val^*(\game, \strategy)
\end{equation*}
over all tensor product strategies.  A strategy is projective if all its
measurements are projective, and symmetric if the state is
permutation-invariant on $\mathcal H \ot \mathcal H$ and both players use the
same measurements.

The symmetrization lemma of \cite[Section~5]{Ji2020MIPStar} asserts:%
\footnote{The statement is at
\path{references/qpbt-paper/06_nonlocal_games_and_mipstar.tex}, lines~94--99,
and its proof at lines~101--132.  The blueprint statement kept in this form
is \path{lem:symmetric-strat} in
\path{blueprint/src/chapter/ch12_qpbt_games.tex}.}
if $\game$ is a symmetric game with $\val^*(\game) = 1-\eps$ for some
$\eps \ge 0$, then
\begin{equation*}
  \exists\, \strategy = (\ket{\psi}, M) \text{ symmetric and projective}
  \quad\text{with}\quad
  \val^*(\game, \strategy) \ge 1-\eps.
  \tag{\path{lem:symmetric-strat}}
\end{equation*}

Its proof runs as follows.  For every $\eps' > \eps$ the definition of the
supremum supplies a strategy $\strategy' = (\ket{\psi'}, A, B)$ with
$\val^*(\game, \strategy') \ge 1-\eps'$.  By Naimark dilation in the completed
projective form of \cite[Theorem~5.1]{Ji2020LowIndividualDegree} one may take
$\ket{\psi'} \in \C^d_{A'} \ot \C^d_{B'}$ for some integer $d$ with all
measurements projective.  One then forms the permutation-invariant state
\begin{equation*}
  \ket{\psi} = \tfrac{1}{\sqrt 2}
  \bigl( \ket{0}_A \ket{1}_B \ket{\psi'}_{A'B'}
       + \ket{1}_A \ket{0}_B \ket{\psi'_\tau}_{A'B'} \bigr)
\end{equation*}
on $(\C^2 \ot \C^d) \ot (\C^2 \ot \C^d)$, where $\ket{\psi'_\tau}$ permutes
the two players' registers, and the block-diagonal projective measurements
$M^x_a = \ket{0}\!\bra{0} \ot A^x_a + \ket{1}\!\bra{1} \ot B^x_a$.  Symmetry
of $D$ gives $\val^*(\game, (\ket\psi, M)) = \val^*(\game, \strategy')$.  The
construction is therefore \emph{value-preserving}: applied to any strategy it
produces a symmetric projective strategy of exactly the same value.  What the
proof establishes is thus the given-strategy form, hence the approximate form;
the printed conclusion is the attainment form.

\section{Why attainment does not follow}

The supremum defining $\val^*(\game)$ ranges over strategies on
finite-dimensional Hilbert spaces of unbounded dimension.  The strategies
$\strategy'_{\eps'}$ produced for $\eps' \downarrow \eps$ may require
dimensions $d(\eps')$ growing without bound, and there is no compactness
across dimensions from which a limiting finite-dimensional strategy of value
at least $1-\eps$ could be extracted.  A symmetric projective strategy of
value at least $1-\eps = \val^*(\game)$ would attain the supremum, so the
printed conclusion is equivalent, under its hypothesis, to the assertion that
$\val^*(\game)$ is attained by some strategy.  No general principle supplies
this: it is known, and recalled in the introduction of the source itself, that
the set of finite-dimensional quantum correlations is not closed and differs
from the set of infinite-dimensional spatial correlations, so limits of
finite-dimensional strategies need not be realized by any finite-dimensional
strategy.  Whether some symmetric game realizes the failure --- a game whose
value is attained by no tensor product strategy, presented together with the
exposing functional as a decision predicate --- is a separate question that
this note does not need to settle.  The local verdict is about the proof: the
cited argument proves the approximate form and does not prove the attainment
form.

\section{An additional obstruction in the formal domain}

The current Lean structure for a symmetric game assumes finiteness of its
answer alphabet but does not require that alphabet to be nonempty.  Consider
the game with question alphabet $\mathcal X=\{*\}$, answer alphabet
$\mathcal A=\varnothing$, the probability distribution concentrated on
$(*,*)$, and the unique decision predicate on the empty answer domain.
Both symmetry requirements hold.  There is no tensor product strategy for
this game: a unit vector in $\mathcal H_A\ot\mathcal H_B$ forces both local
spaces to be nonzero, whereas completeness of Alice's empty measurement
would give
\[
  0=\sum_{a\in\varnothing} A^*_a=I_{\mathcal H_A},
\]
which is impossible on a nonzero space.

Lean defines the game value using the conditional supremum operation on
$\mathbb R$.  For the empty set this operation is assigned the value zero,
as expressed by Mathlib's \leanid{Real.sSup\_empty}.  This is a convention
outside the nonempty domain on which the usual least-upper-bound property
holds; it is not an attainment principle.  Thus \leanid{Game.value} assigns
the example the value $0=1-\eps$ with $\eps=1\geq0$, but it has no symmetric
projective strategy, since it has no strategy at all.  Consequently the exact
current Lean statement is false on its admitted domain, independently of the
unbounded-dimension issue.  The complete verification and its axiom audit are
recorded in \path{audits/2026-09-12_issue-524-symmetrization-obstruction.md}.%
\footnote{The verification reuses the existing
\leanid{nonempty\_of\_unit\_vector} and
\leanid{measurement\_outcome\_nonempty} lemmas.  Its declarations depend only
on \leanid{propext}, \leanid{Classical.choice}, and \leanid{Quot.sound},
not on the unfinished symmetrization theorem.}

This example concerns a boundary of the formal encoding.  Although the
printed definition of a game lists finite alphabets without an explicit
nonemptiness condition, it does not specify a real value for the supremum
over an empty strategy class.  The example therefore does not establish
nonattainment for the intended class of games with nonempty alphabets.
Making that domain condition explicit would exclude this example but would
leave the source proof's attainment step unproved.  The statement-integrity
review requested here was carried out on 2026-09-19 and is the statement
correction recorded below; the printed Lean statement and its proof-gap marker
were replaced there by an unasserted Lean proposition and this refutation,
proved as \leanid{not\_forall\_printedSymmetricProjectiveAttainmentClaim}.

\section{Uses of the lemma}

The source invokes the lemma twice.

First, in the soundness proof for the oracularization transformation, the
argument fixes a strategy $\strategy^{\mathrm{orac}}$ for the oracularized game
``with value $1-\eps$'' and applies the lemma to assume without loss of
generality that it is projective.%
\footnote{\path{references/qpbt-paper/10_oracularization.tex},
lines~317--321.}
Here the input is a given strategy, not the value of the game; the
given-strategy form applies verbatim and preserves the value exactly, so the
use is justified by what the proof establishes.

Second, in the soundness proof for answer reduction, the argument assumes the
strict inequality $\val^*(\game) > 1-\eps$ and concludes that there is a
symmetric projective strategy of value greater than $1-\eps$.%
\footnote{\path{references/qpbt-paper/11_answer_reduction.tex},
lines~2322--2332.}
Writing $\val^*(\game) = 1-\eps_0$ with $\eps_0 < \eps$ and applying the
approximate form with $\eps' = (\eps+\eps_0)/2$ produces a symmetric
projective strategy of value at least $1-\eps' > 1-\eps$.  A strict lower
bound on the supremum thus passes through the approximate form; attainment is
not used.

In this project, the soundness statements of the Pauli basis test track
quantify over strategies assumed to succeed with probability at least
$1-\eps$, so any symmetrization step starts from a given strategy and consumes
the given-strategy form.  The blueprint records both forms as part of the
games vocabulary and has no downstream dependency on either; the two source
uses above belong to transformations outside the scope of the present
development.
Consequently, no consumer of the lemma --- in the source or here --- requires
the attainment form.

\section{Comparison with the blueprint}

The two subsections that follow describe the arrangement in force before
2026-09-19; the statement-correction section below supersedes them.

The blueprint \cite{MIPStarREBlueprint} states the lemma in the source shape,
as the proof-gap protocol requires when the printed statement is not known to
be false and only its proof is blocked.  Its proof records the
value-preserving construction and the approximate form that the construction
yields, and explicitly marks the remaining attainment step as unproved, citing
this note; an accompanying remark summarizes the obstruction.  The
value-preserving given-strategy form is stated beside it as a separate
formalization-support node, \path{lem:symmetric-strat-given-strategy}, which
is not advertised as the source lemma.

The Lean declarations follow the same separation; the realignment section
of 2026-09-18 below records two further ones.  The source-shaped
\leanid{exists\_symmetric\_projective\_strategy} is linked to
\path{lem:symmetric-strat}; its proof is the tracked attainment obligation and
remains open.  The given-strategy
\leanid{exists\_symmetric\_projective\_strategy\_of\_strategy}, linked to the
support node, has the same hypothesis, exact value equality, and projective
symmetric conclusion as that node, and its proof is complete.

\section{Realignment of the Lean statements (2026-09-18)}

This realignment is published as \ghpr{601}.  Following the disposition
recorded for \ghissue{524}, the formalization keeps
the source-labelled statement as printed and adds the approximate form beside
it as a separate, proved Lean declaration.  Three points fix the record.

First, the printed attainment form remains an \emph{open source gap}.  The
declaration \leanid{exists\_symmetric\_projective\_strategy} keeps the source
shape, keeps its proof hole, and keeps its unfaithfulness marker; the blueprint
node \path{lem:symmetric-strat} is still stated as printed and still carries no
proof-level \texttt{leanok} flag.  Nothing in this realignment proves the
printed statement, and the weaker result is not renamed as it.

Second, the approximate form listed under ``Key theorem forms'' above is now
formalized and proved as
\leanid{exists\_symmetric\_projective\_strategy\_approx}, with the equivalent
slack form \leanid{exists\_symmetric\_projective\_strategy\_of\_lt\_value}
beside it, against a new formalization-only blueprint node
\path{lem:symmetric-strat-approx}.  The Lean proof is the source argument
unchanged: the supremum property of $\val^*(\game)$ supplies a strategy above
the strict lower bound, and the already proved value-preserving construction
\leanid{exists\_symmetric\_projective\_strategy\_of\_strategy} symmetrizes it.
The auxiliary facts added for it --- that every strategy value is at most one,
that the set of strategy values is bounded above, and that a nonempty answer
alphabet admits a strategy --- are formalization-only support for
\path{def:tensor-product-value}.

Third, the nonempty answer alphabet appears as an explicit hypothesis of both
approximate declarations.  This is the boundary condition of the formal-domain
obstruction: with an empty answer alphabet the Lean game has no strategy at
all, while $\val^*(\game)$ is still defined, as the totalized supremum of the
empty set, so the statement as printed is false there.  That obstruction and
its checked Lean witness are the subject of the separate documentation packet
for \ghissue{524}; the present section does not restate or re-audit its
evidence.

\section{Statement correction (2026-09-19)}

This section is published as the zero-sorry statement correction for
\ghissue{524}, on the owner decision of 2026-09-19 that the development carry
no \texttt{sorry}.  It supersedes the disposition of the realignment section
above, which kept the printed form as an unproved Lean theorem.

The reclassification anticipated in the conclusion below has now happened, and
for the reason recorded in the formal-domain section: on the domain over which
the Lean statement was quantified, the printed form is not merely unproved, it
is \emph{false}.  Keeping it as a Lean theorem with an open proof therefore
asserted something the development cannot establish and, on part of its own
domain, could never establish.  Four changes implement the correction; none of
them weakens a statement that was proved, and none of them removes the
discrepancy from view.

\begin{enumerate}
  \item The Lean theorem \leanid{exists\_symmetric\_projective\_strategy},
  with its \texttt{sorry}, is removed.  Nothing in the development depended on
  it.  Its printed content is preserved verbatim, and unasserted, as the
  Prop-valued definition
  \leanid{PrintedSymmetricProjectiveAttainmentClaim}, whose binders are those
  of the printed sentence and whose docstring records the locator, the reason
  the sentence is not established, and the pointers to the proved forms and to
  this note.
  \item The refutation recorded in the formal-domain section above is now a
  proved Lean theorem rather than an audit artifact:
  \leanid{not\_forall\_printedSymmetricProjectiveAttainmentClaim} proves that
  the printed claim does not hold for every symmetric game and every $\eps$,
  using the empty-answer game
  \leanid{SymmetrizationObstruction.emptyAnswerGame} together with
  \leanid{SymmetrizationObstruction.no\_strategy} and
  \leanid{SymmetrizationObstruction.value\_eq\_zero}.  These are the
  declarations checked in
  \path{audits/2026-09-12_issue-524-symmetrization-obstruction.md}, moved into
  the library unchanged.
  \item The blueprint node \path{lem:symmetric-strat} now carries the
  \emph{corrected} statement --- the approximate form, with the nonempty
  answer alphabet displayed --- with \leanid{exists\_symmetric\_projective\_strategy\_approx}
  and \leanid{exists\_symmetric\_projective\_strategy\_of\_lt\_value} as its
  formalization and with statement- and proof-level \texttt{leanok}.  The
  printed sentence is stated in full in \path{rem:symmetric-strat-limit}
  beside it, and the Lean record and its refutation are the node
  \path{lem:symmetric-strat-printed-claim}.  The former approximate node
  \path{lem:symmetric-strat-approx} is folded into \path{lem:symmetric-strat}
  and no longer exists.
  \item The given-strategy node \path{lem:symmetric-strat-given-strategy} and
  its Lean declaration are unchanged.
\end{enumerate}

This is a correction of the \emph{statement}, in the sense of the register:
the mathematics the source argument establishes is stated and proved, and the
mathematics the source prints is displayed beside it and shown not to follow.
The approximate assertion is weaker than the printed attainment assertion,
even though this change removes no previously proved result. The unasserted
proposition and the empty-answer refutation do not establish attainment on the
source's nonempty-answer domain. That distinction remains relevant to the
adoption conditions of \path{local/protocols/completion.md}; the change alone
does not give this row a terminal status.

\section{Conclusion}

For games with nonempty answer alphabets, the verdict remains a proof-level
gap in the cited lemma: the symmetrization argument proves the given-strategy
and approximate forms, while the printed
conclusion additionally asserts that the supremum defining $\val^*(\game)$ is
attained, which does not follow.  The gap is harmless for the results built
on the lemma, since both uses in the source, and all planned uses in this
project, proceed from a given strategy of value at least $1-\eps$ or from a
strict lower bound on the value, for which the value-preserving construction
suffices.  Stage~4 has carried out the recommendation of this note: the
given-strategy symmetrization lemma is formalized and proved as the workhorse
statement.  The 2026-09-18
realignment added the approximate form as a proved Lean declaration under an
explicit nonempty-answer hypothesis, and the 2026-09-19 statement correction
above makes that approximate form the formalization of the source node, keeps
the printed sentence visible as an unasserted Lean proposition, and proves its
refutation on the formal domain.  The reclassification from an unproved step to
a documented statement correction, anticipated here, is thereby carried out.

The exact current Lean domain has the empty-answer counterexample
above.  An unrestricted attainment theorem cannot discharge the printed
statement.  Neither obstruction is discharged by the correction; both are made
explicit by it.

\bibliographystyle{alpha}
\bibliography{references}

\end{document}
